\documentclass[11pt]{article}

\usepackage[utf8]{inputenc}
\usepackage[T1]{fontenc}
\usepackage{lmodern}
\usepackage{geometry}
\usepackage{amsmath,amssymb,amsfonts,amsthm,mathtools}
\usepackage{bm}
\usepackage{enumitem}
\usepackage{microtype}
\usepackage{hyperref}
\usepackage{titlesec}
\usepackage{booktabs}
\usepackage{array}
\usepackage{setspace}
\usepackage{mathrsfs}
\usepackage{physics}

\geometry{margin=1in}
\hypersetup{colorlinks=true, linkcolor=black, urlcolor=blue, citecolor=black}
\setlist[enumerate]{topsep=0.4em,itemsep=0.35em}
\setlist[itemize]{topsep=0.3em,itemsep=0.25em}
\titleformat{\section}{\large\bfseries}{\thesection.}{0.5em}{}
\titleformat{\subsection}{\normalsize\bfseries}{\thesubsection.}{0.5em}{}
\setstretch{1.06}

\newcommand{\Aether}{\AE{}ther}
\newcommand{\Aflow}{\AE{}ther-flow}
\newcommand{\TheoryName}{\Aflow{} Interpretation of Relativity}
\newcommand{\TheoryProgram}{\Aether{} / \Aflow{} framework}

\newcommand{\CompletedMetric}{g^{\sharp}}
\newcommand{\MatterSpace}{\Psi}

\newcommand{\Car}{\mathrm{car}}
\newcommand{\Cur}{\mathrm{cur}}
\newcommand{\Hom}{\mathrm{hom}}

\newcommand{\ForSpace}[1]{\mathcal I_{#1}^{\mathrm{for}}}
\newcommand{\PrimShell}{\mathcal S_{\mathrm{prim}}}
\newcommand{\Reduction}[1]{\mathcal R_{#1}}
\newcommand{\CurrentCarrier}[1]{\mathfrak C_{#1}^{\mathrm{cur}}}
\newcommand{\EqCarrier}[1]{\mathfrak C_{#1}^{\mathrm{eq}}}
\newcommand{\MetricOut}{\mathscr G_{\mathrm{out}}}
\newcommand{\MatterOut}{\mathscr T_{\mathrm{out}}}

\newcommand{\FoundSpace}[1]{\mathcal F_{#1}^{\mathrm{fdn}}}
\newcommand{\GroundSpace}[1]{\mathcal G_{#1}^{\mathrm{grd}}}
\newcommand{\OrigSpace}[1]{\mathcal O_{#1}^{\mathrm{orig}}}
\newcommand{\PrecSpace}[1]{\mathcal P_{#1}^{\mathrm{prec}}}
\newcommand{\LiftSpace}[1]{\mathcal L_{#1}^{\mathrm{lift}}}
\newcommand{\SubSpace}[1]{\mathcal M_{#1}^{\mathrm{sub}}}
\newcommand{\SubTarget}[1]{E_{#1}^{\mathrm{sub}}}
\newcommand{\SubPair}[1]{\mathfrak s_{#1}^{\mathrm{sub}}}
\newcommand{\EqnTarget}[1]{Q_{#1}^{\mathrm{eqn}}}

\newcommand{\HiddenSpace}[1]{\mathcal H_{#1}^{\mathrm{hid}}}
\newcommand{\HiddenBody}[1]{D_{#1}^{\mathrm{hid}}}

\newcommand{\SeedPackage}{\mathfrak S^{\mathrm{seed}}}
\newcommand{\SeedSpace}[1]{\mathcal S_{#1}^{\mathrm{seed}}}
\newcommand{\SeedBody}[1]{\mathcal B_{#1}^{\mathrm{seed}}}
\newcommand{\SeedProj}[1]{\Pi_{#1}^{\mathrm{seed}}}
\newcommand{\SeedFoundMin}[1]{\mathfrak f_{#1}^{\mathrm{seed,fdn}}}
\newcommand{\SeedGroundMin}[1]{\mathfrak f_{#1}^{\mathrm{seed,grd}}}
\newcommand{\SeedOrigMin}[1]{\mathfrak f_{#1}^{\mathrm{seed,orig}}}
\newcommand{\SeedPrecMin}[1]{\mathfrak f_{#1}^{\mathrm{seed,prec}}}
\newcommand{\SeedLiftMin}[1]{\mathfrak f_{#1}^{\mathrm{seed,lift}}}
\newcommand{\SeedSubMin}[1]{\mathfrak m_{#1}^{\mathrm{seed}}}
\newcommand{\SeedSection}[1]{\Gamma_{#1}^{\mathrm{seed}}}
\newcommand{\SeedFunctional}[1]{\mathscr W_{#1}^{\mathrm{seed}}}
\newcommand{\SeedShellSet}[1]{\mathcal Z_{#1}^{\mathrm{seed}}}
\newcommand{\SeedZeroCore}[1]{\mathcal Z_{#1}^{0,\mathrm{seed}}}
\newcommand{\SeedSplit}[1]{\Phi_{#1}^{\mathrm{seed}}}
\newcommand{\SeedCharge}[1]{\mathcal Q_{#1}^{\mathrm{seed}}}
\newcommand{\SeedEmbed}[1]{\zeta_{#1}^{\mathrm{seed}}}
\newcommand{\SeedKernelCh}[1]{\widetilde K_{#1}^{0,\mathrm{seed}}}
\newcommand{\SeedStateComp}[1]{P_{#1}^{\mathrm{seed,co}}}

\newcommand{\GermPackage}{\mathfrak G^{\mathrm{germ}}}
\newcommand{\GermSpace}[1]{\mathcal G_{#1}^{\mathrm{germ}}}
\newcommand{\GermBody}[1]{\mathcal B_{#1}^{\mathrm{germ}}}
\newcommand{\GermProj}[1]{\Pi_{#1}^{\mathrm{germ}}}
\newcommand{\GermSeedMin}[1]{\mathfrak f_{#1}^{\mathrm{germ,seed}}}
\newcommand{\GermFoundMin}[1]{\mathfrak f_{#1}^{\mathrm{germ,fdn}}}
\newcommand{\GermGroundMin}[1]{\mathfrak f_{#1}^{\mathrm{germ,grd}}}
\newcommand{\GermOrigMin}[1]{\mathfrak f_{#1}^{\mathrm{germ,orig}}}
\newcommand{\GermPrecMin}[1]{\mathfrak f_{#1}^{\mathrm{germ,prec}}}
\newcommand{\GermLiftMin}[1]{\mathfrak f_{#1}^{\mathrm{germ,lift}}}
\newcommand{\GermSubMin}[1]{\mathfrak m_{#1}^{\mathrm{germ}}}
\newcommand{\GermSection}[1]{\Gamma_{#1}^{\mathrm{germ}}}
\newcommand{\GermFunctional}[1]{\mathscr W_{#1}^{\mathrm{germ}}}
\newcommand{\GermShellSet}[1]{\mathcal Z_{#1}^{\mathrm{germ}}}
\newcommand{\GermZeroCore}[1]{\mathcal Z_{#1}^{0,\mathrm{germ}}}
\newcommand{\GermSplit}[1]{\Phi_{#1}^{\mathrm{germ}}}
\newcommand{\GermCharge}[1]{\mathcal Q_{#1}^{\mathrm{germ}}}
\newcommand{\GermEmbed}[1]{\zeta_{#1}^{\mathrm{germ}}}
\newcommand{\GermKernelCh}[1]{\widetilde K_{#1}^{0,\mathrm{germ}}}
\newcommand{\GermStateComp}[1]{P_{#1}^{\mathrm{germ,co}}}

\newcommand{\ProtoPackage}{\mathfrak P^{\mathrm{proto}}}
\newcommand{\ProtoSpace}[1]{\mathcal U_{#1}^{\mathrm{proto}}}
\newcommand{\ProtoBody}[1]{\mathcal B_{#1}^{\mathrm{proto}}}
\newcommand{\ProtoProj}[1]{\Pi_{#1}^{\mathrm{proto}}}
\newcommand{\ProtoGermMin}[1]{\mathfrak f_{#1}^{\mathrm{proto,germ}}}
\newcommand{\ProtoSeedMin}[1]{\mathfrak f_{#1}^{\mathrm{proto,seed}}}
\newcommand{\ProtoFoundMin}[1]{\mathfrak f_{#1}^{\mathrm{proto,fdn}}}
\newcommand{\ProtoGroundMin}[1]{\mathfrak f_{#1}^{\mathrm{proto,grd}}}
\newcommand{\ProtoOrigMin}[1]{\mathfrak f_{#1}^{\mathrm{proto,orig}}}
\newcommand{\ProtoPrecMin}[1]{\mathfrak f_{#1}^{\mathrm{proto,prec}}}
\newcommand{\ProtoLiftMin}[1]{\mathfrak f_{#1}^{\mathrm{proto,lift}}}
\newcommand{\ProtoSubMin}[1]{\mathfrak m_{#1}^{\mathrm{proto}}}
\newcommand{\ProtoSection}[1]{\Gamma_{#1}^{\mathrm{proto}}}
\newcommand{\ProtoFunctional}[1]{\mathscr W_{#1}^{\mathrm{proto}}}
\newcommand{\ProtoShellSet}[1]{\mathcal Z_{#1}^{\mathrm{proto}}}
\newcommand{\ProtoZeroCore}[1]{\mathcal Z_{#1}^{0,\mathrm{proto}}}
\newcommand{\ProtoSplit}[1]{\Phi_{#1}^{\mathrm{proto}}}
\newcommand{\ProtoCharge}[1]{\mathcal Q_{#1}^{\mathrm{proto}}}
\newcommand{\ProtoEmbed}[1]{\zeta_{#1}^{\mathrm{proto}}}
\newcommand{\ProtoKernelCh}[1]{\widetilde K_{#1}^{0,\mathrm{proto}}}
\newcommand{\ProtoStateComp}[1]{P_{#1}^{\mathrm{proto,co}}}

\newtheorem{definition}{Definition}
\newtheorem{proposition}{Proposition}
\newtheorem{remark}{Remark}
\newtheorem{corollary}{Corollary}
\newtheorem{theorem}{Theorem}

\title{\textbf{The \TheoryName{}}\\[0.4em]
\large Primitive-Reservoir Observer-Reduction\\
\large Seed-Generating Substrate Germ Dynamics\\
\large from Germ-Generating Substrate Proto-Germ Dynamics}
\author{}
\date{}

\begin{document}

\maketitle

\begin{abstract}
The active primitive-reservoir observer-reduction frontier is now the theorem deriving the current seed package from seed-generating substrate germ dynamics. Under the recorded routing safeguard, the honest next benchmark-facing move must therefore do one of two things: either derive or justify the germ package from still more primitive substrate structure, or prove a genuinely smaller theorem-level collapse strictly inside the observer-relevant germ package. The present note takes the default first route.

For each sector it introduces \emph{germ-generating substrate proto-germ dynamics}: a compact convex proto-germ body over the recorded germ state space, unique proto-germ minimizers on the fixed germ, seed, foundation, ground, origin, precursor, lift, and substrate fibers, a smooth proto-germ restriction section, an exact proto-germ shell projecting diffeomorphically to the exact germ shell, a closed embedded proto-germ zero core whose projected split reproduces the germ split, a hidden charge, an observer-compatible exact proto-germ zero-response realization, and a fixed-data solvability / coercive package projecting to the recorded germ kernel data. The current seed-generating substrate germ dynamics are then recovered canonically as the projected exact-shell transport image of that deeper proto-germ package.

The benchmark boundary remains fixed throughout. The one operative metric is still exactly \(\CompletedMetric\), the ordinary matter route is unchanged, there is no second operative metric, and the low-energy matter law is not enlarged. The positive result is therefore narrow but benchmark facing: the current germ package is no longer left as an unexplained primitive stopping point on the active line. The remaining honest burden moves one layer deeper again, to the proto-germ dynamics themselves or to a sharper collapse inside their observer-relevant core.
\end{abstract}

\section{Introduction}

The active derivational program of \TheoryName{} has now reached a sharper burden than another same-output germ restatement. The current germ theorem already showed that the present seed package is the exact-shell transport image of a more primitive seed-generating substrate germ dynamics package. Once that result is read together with the benchmark-routing safeguard against recursive depth drift, the next honest move becomes narrower: either derive or justify the germ package from still more primitive substrate structure, or prove a genuinely smaller theorem-level collapse strictly inside the observer-relevant germ package. Reopening the older root-nucleus handoff, re-running another same-output seed / root / foundation relay, or returning to stationary-reduction, stationary-Hessian, spectral, or selector layers would no longer address the live burden.

The present note takes the first route. Its positive claim is not that final substrate microphysics has been identified, and it is not that the benchmark exact-relativistic package has been changed. Its claim is narrower and benchmark facing. The current germ package is shown to arise canonically from a more primitive proto-germ-level substrate dynamics whose projected exact-shell transport already contains the observer-relevant structure of the germ package. In that sense the germ layer is no longer the primitive place where the active line has to stop.

Two features matter here. First, the theorem targets the current germ package itself rather than revisiting the already-spent seed handoff, so it does not spend the next move on an older same-output presentation. Second, the deeper proto-germ package is not merely another renaming of the germ package. It adds an explicit proto-germ state space, proto-germ body, proto-to-germ projection, and deeper exact-shell transport data below the germ frontier. The germ shell, germ zero core, split, and lower minimizer hierarchy are therefore projected images of a deeper substrate package rather than unexplained primitive data.

\paragraph{Framework and claim boundary.}
\TheoryProgram{} denotes the broader \Aether{} / \Aflow{} framework built on the claims that the \Aether{} is the underlying four-dimensional substrate of reality, the \Aflow{} is its intrinsic ordered motion, observed three-dimensional space is the local experiential slice of that deeper substrate, S-time is the experienced order of change arising from matter, light, and the \Aflow{}, and observed expansion is the three-dimensional appearance of deeper four-dimensional ordered motion. Within that framework, \TheoryName{} denotes the exact-closure theory in which the effective gravitational dynamics are adopted to be exactly Einsteinian with universal matter coupling. In the active sequence, \emph{adoption} means use of that established relativistic dynamics without claiming substrate derivation, while \emph{derivation} is reserved for a first-principles recovery from explicit substrate variables.

The present note stays strictly inside that boundary. It does not alter the benchmark exact-closure package. It does not introduce a second operative metric or a new low-energy matter law. It only proves that the current seed-generating substrate germ dynamics can itself be generated from a more primitive proto-germ-level substrate dynamics.

\section{Fixed Inputs}

\begin{definition}[Recorded primitive continuation data]
\label{def:fixed}
Fix the following continuation data from the active primitive-reservoir observer-reduction line.
\begin{enumerate}
    \item For each sector label \(\sigma\in\{\Car,\Cur,\Hom\}\), a primitive forcing shell
    \[
    \ForSpace{\sigma},
    \]
    a visible reduction map
    \[
    \Reduction{\sigma}:\ForSpace{\sigma}\to X_{\sigma},
    \]
    and shell-level current and equilibrium carriers
    \[
    \CurrentCarrier{\sigma}:\ForSpace{\sigma}\to Z_{\sigma}^{\mathrm{cur}},
    \qquad
    \EqCarrier{\sigma}:\ForSpace{\sigma}\to Z_{\sigma}^{\mathrm{eq}}.
    \]
    \item The product shell
    \[
    \PrimShell
    \equiv
    \ForSpace{\Car}\times \ForSpace{\Cur}\times \ForSpace{\Hom}.
    \]
    \item The recorded metric and ordinary-matter outputs
    \[
    \MetricOut:\PrimShell\to\{\text{observer metrics}\},
    \qquad
    \MatterOut:\PrimShell\times\MatterSpace\to\{\text{observer stress tensors}\},
    \]
    satisfying
    \begin{equation}
    \MetricOut(\mathbf w)=\CompletedMetric,
    \qquad
    \MatterOut(\mathbf w,\psi)
    =
    T_{\mu\nu}^{\mathrm{matter}}[\CompletedMetric,\psi]
    \label{eq:fixedoutputs}
    \end{equation}
    for every \(\mathbf w\in\PrimShell\) and every matter configuration \(\psi\in\MatterSpace\).
\end{enumerate}
\end{definition}

\begin{definition}[Recorded lower-fiber ladder]
\label{def:ladder}
Fix \(\sigma\in\{\Car,\Cur,\Hom\}\). The active line already fixes the lower-fiber ladder
\[
\FoundSpace{\sigma},
\qquad
\GroundSpace{\sigma},
\qquad
\OrigSpace{\sigma},
\qquad
\PrecSpace{\sigma},
\qquad
\LiftSpace{\sigma},
\qquad
\SubSpace{\sigma},
\]
together with the common substrate target \(\SubTarget{\sigma}\), the positive-definite pairing
\[
\SubPair{\sigma}:
\SubTarget{\sigma}\times\SubTarget{\sigma}\to\mathbb R,
\]
and the hidden equation target \(\EqnTarget{\sigma}\).
\end{definition}

\begin{definition}[Recorded downstream seed package]
\label{def:seed}
Fix \(\sigma\in\{\Car,\Cur,\Hom\}\). The active line already carries a recorded seed package \(\SeedPackage\), consisting of:
\begin{enumerate}
    \item a seed state space \(\SeedSpace{\sigma}\), compact convex seed body \(\SeedBody{\sigma}\), affine seed projection \(\SeedProj{\sigma}:\SeedSpace{\sigma}\to\FoundSpace{\sigma}\), and the fixed seed minimizers
    \[
    \SeedFoundMin{\sigma},\quad
    \SeedGroundMin{\sigma},\quad
    \SeedOrigMin{\sigma},\quad
    \SeedPrecMin{\sigma},\quad
    \SeedLiftMin{\sigma},\quad
    \SeedSubMin{\sigma};
    \]
    \item a seed restriction section \(\SeedSection{\sigma}\), seed functional \(\SeedFunctional{\sigma}\), and exact seed shell \(\SeedShellSet{\sigma}\);
    \item a closed embedded seed zero core \(\SeedZeroCore{\sigma}\), hidden fiber space \(\HiddenSpace{\sigma}\) with compact hidden body \(\HiddenBody{\sigma}\), and hidden split diffeomorphism \(\SeedSplit{\sigma}\);
    \item the hidden charge \(\SeedCharge{\sigma}\), exact seed zero-response realization \(\SeedEmbed{\sigma}\), and fixed-data solvability / coercive package recorded by \(\SeedKernelCh{\sigma}\) and \(\SeedStateComp{\sigma}\).
\end{enumerate}
\end{definition}

\begin{remark}
Definitions \ref{def:fixed}--\ref{def:seed} are not reopened here. The primitive continuation data, the one operative metric, the ordinary matter route, the recorded lower-fiber ladder, and the downstream seed package remain fixed inputs. The present note addresses only whether the current germ package can itself be generated from more primitive substrate dynamics.
\end{remark}

\section{Recorded Germ Target}

\begin{definition}[Recorded seed-generating substrate germ dynamics]
\label{def:germ}
Fix \(\sigma\in\{\Car,\Cur,\Hom\}\). The recorded germ package \(\GermPackage\) for sector \(\sigma\) consists of the following data.
\begin{enumerate}
    \item A finite-dimensional Euclidean germ state space \(\GermSpace{\sigma}\), a compact convex germ body
    \[
    \GermBody{\sigma}\subset\GermSpace{\sigma}
    \]
    with nonempty interior, an affine surjection
    \[
    \GermProj{\sigma}:\GermSpace{\sigma}\to\SeedSpace{\sigma},
    \]
    and smooth germ fiber minimizers
    \[
    \GermSeedMin{\sigma},\quad
    \GermFoundMin{\sigma},\quad
    \GermGroundMin{\sigma},\quad
    \GermOrigMin{\sigma},\quad
    \GermPrecMin{\sigma},\quad
    \GermLiftMin{\sigma},\quad
    \GermSubMin{\sigma}
    \]
    satisfying
    \[
    \GermProj{\sigma}\bigl(\GermBody{\sigma}\bigr)=\SeedBody{\sigma},
    \qquad
    \GermProj{\sigma}\circ\GermSeedMin{\sigma}
    =
    \mathrm{id}_{\SeedBody{\sigma}},
    \]
    together with the transport identities
    \[
    \GermProj{\sigma}\circ\GermFoundMin{\sigma}
    =
    \SeedFoundMin{\sigma},
    \qquad
    \GermProj{\sigma}\circ\GermGroundMin{\sigma}
    =
    \SeedGroundMin{\sigma},
    \]
    \[
    \GermProj{\sigma}\circ\GermOrigMin{\sigma}
    =
    \SeedOrigMin{\sigma},
    \qquad
    \GermProj{\sigma}\circ\GermPrecMin{\sigma}
    =
    \SeedPrecMin{\sigma},
    \]
    \[
    \GermProj{\sigma}\circ\GermLiftMin{\sigma}
    =
    \SeedLiftMin{\sigma},
    \qquad
    \GermProj{\sigma}\circ\GermSubMin{\sigma}
    =
    \SeedSubMin{\sigma}.
    \]
    \item A smooth germ restriction section
    \[
    \GermSection{\sigma}:\GermSpace{\sigma}\to\SubTarget{\sigma},
    \]
    the germ functional
    \[
    \GermFunctional{\sigma}(g)
    \equiv
    \SubPair{\sigma}\bigl(\GermSection{\sigma}(g),\GermSection{\sigma}(g)\bigr),
    \]
    and the exact germ shell
    \[
    \GermShellSet{\sigma}
    \equiv
    \bigl(\GermSection{\sigma}\bigr)^{-1}(0)\cap\GermBody{\sigma},
    \]
    satisfying
    \[
    \GermSection{\sigma}\circ\GermSeedMin{\sigma}
    =
    \SeedSection{\sigma},
    \qquad
    \GermProj{\sigma}\big|_{\GermShellSet{\sigma}}
    :
    \GermShellSet{\sigma}
    \xrightarrow{\cong}
    \SeedShellSet{\sigma}.
    \]
    \item A closed embedded germ zero core
    \[
    \GermZeroCore{\sigma}\subset\GermShellSet{\sigma},
    \]
    the same hidden fiber space \(\HiddenSpace{\sigma}\) with the same compact hidden body \(\HiddenBody{\sigma}\subset\HiddenSpace{\sigma}\), a diffeomorphism
    \[
    \GermProj{\sigma}\big|_{\GermZeroCore{\sigma}}
    :
    \GermZeroCore{\sigma}
    \xrightarrow{\cong}
    \SeedZeroCore{\sigma},
    \]
    and a split diffeomorphism
    \[
    \GermSplit{\sigma}:
    \GermZeroCore{\sigma}\times\HiddenBody{\sigma}
    \xrightarrow{\cong}
    \GermShellSet{\sigma}
    \]
    obeying the projected split identity
    \[
    \SeedSplit{\sigma}(z,\eta)
    =
    \GermProj{\sigma}\Bigl(
    \GermSplit{\sigma}\bigl(
    (\GermProj{\sigma}\big|_{\GermZeroCore{\sigma}})^{-1}(z),\eta
    \bigr)\Bigr)
    \]
    for every \(z\in\SeedZeroCore{\sigma}\) and every \(\eta\in\HiddenBody{\sigma}\).
    \item A hidden charge
    \[
    \GermCharge{\sigma}:\HiddenSpace{\sigma}\to\EqnTarget{\sigma},
    \]
    a smooth observer-compatible exact germ zero-response realization
    \[
    \GermEmbed{\sigma}:\ForSpace{\sigma}\hookrightarrow\GermZeroCore{\sigma},
    \]
    satisfying
    \[
    \SeedCharge{\sigma}
    =
    \GermCharge{\sigma},
    \qquad
    \SeedEmbed{\sigma}
    =
    \GermProj{\sigma}\circ\GermEmbed{\sigma},
    \]
    and a fixed-data solvability / coercive package on \(\GermZeroCore{\sigma}\), recorded through \(\GermKernelCh{\sigma}\) and \(\GermStateComp{\sigma}\), whose projection through
    \[
    \GermProj{\sigma}\big|_{\GermZeroCore{\sigma}}
    \]
    recovers the recorded seed kernel chart \(\SeedKernelCh{\sigma}\) and orthogonal projector \(\SeedStateComp{\sigma}\).
\end{enumerate}
\end{definition}

\begin{remark}
Definition \ref{def:germ} is the fixed target already carried by the current germ-frontier theorem. The present note does not spend its move on a second germ-level restatement. It records the current germ package only because the deeper proto-germ dynamics are defined relative to it.
\end{remark}

\section{Germ-Generating Substrate Proto-Germ Dynamics}

\begin{definition}[Germ-generating substrate proto-germ dynamics]
\label{def:proto}
Fix \(\sigma\in\{\Car,\Cur,\Hom\}\). A germ-generating substrate proto-germ dynamics package \(\ProtoPackage\) for sector \(\sigma\) consists of the following data.
\begin{enumerate}
    \item A finite-dimensional Euclidean proto-germ state space \(\ProtoSpace{\sigma}\), a compact convex proto-germ body
    \[
    \ProtoBody{\sigma}\subset\ProtoSpace{\sigma}
    \]
    with nonempty interior, an affine surjection
    \[
    \ProtoProj{\sigma}:\ProtoSpace{\sigma}\to\GermSpace{\sigma},
    \]
    and smooth proto-germ fiber minimizers
    \[
    \ProtoGermMin{\sigma},\quad
    \ProtoSeedMin{\sigma},\quad
    \ProtoFoundMin{\sigma},\quad
    \ProtoGroundMin{\sigma},\quad
    \ProtoOrigMin{\sigma},\quad
    \ProtoPrecMin{\sigma},\quad
    \ProtoLiftMin{\sigma},\quad
    \ProtoSubMin{\sigma}
    \]
    satisfying
    \[
    \ProtoProj{\sigma}\bigl(\ProtoBody{\sigma}\bigr)=\GermBody{\sigma},
    \qquad
    \ProtoProj{\sigma}\circ\ProtoGermMin{\sigma}
    =
    \mathrm{id}_{\GermBody{\sigma}},
    \]
    together with the transport identities
    \[
    \ProtoProj{\sigma}\circ\ProtoSeedMin{\sigma}
    =
    \GermSeedMin{\sigma},
    \qquad
    \ProtoProj{\sigma}\circ\ProtoFoundMin{\sigma}
    =
    \GermFoundMin{\sigma},
    \]
    \[
    \ProtoProj{\sigma}\circ\ProtoGroundMin{\sigma}
    =
    \GermGroundMin{\sigma},
    \qquad
    \ProtoProj{\sigma}\circ\ProtoOrigMin{\sigma}
    =
    \GermOrigMin{\sigma},
    \]
    \[
    \ProtoProj{\sigma}\circ\ProtoPrecMin{\sigma}
    =
    \GermPrecMin{\sigma},
    \qquad
    \ProtoProj{\sigma}\circ\ProtoLiftMin{\sigma}
    =
    \GermLiftMin{\sigma},
    \]
    \[
    \ProtoProj{\sigma}\circ\ProtoSubMin{\sigma}
    =
    \GermSubMin{\sigma}.
    \]
    \item A smooth proto-germ restriction section
    \[
    \ProtoSection{\sigma}:\ProtoSpace{\sigma}\to\SubTarget{\sigma},
    \]
    the proto-germ functional
    \[
    \ProtoFunctional{\sigma}(p)
    \equiv
    \SubPair{\sigma}\bigl(\ProtoSection{\sigma}(p),\ProtoSection{\sigma}(p)\bigr),
    \]
    and the exact proto-germ shell
    \[
    \ProtoShellSet{\sigma}
    \equiv
    \bigl(\ProtoSection{\sigma}\bigr)^{-1}(0)\cap\ProtoBody{\sigma},
    \]
    satisfying
    \[
    \ProtoSection{\sigma}\circ\ProtoGermMin{\sigma}
    =
    \GermSection{\sigma},
    \qquad
    \ProtoProj{\sigma}\big|_{\ProtoShellSet{\sigma}}
    :
    \ProtoShellSet{\sigma}
    \xrightarrow{\cong}
    \GermShellSet{\sigma}.
    \]
    \item A closed embedded proto-germ zero core
    \[
    \ProtoZeroCore{\sigma}\subset\ProtoShellSet{\sigma},
    \]
    the same hidden fiber space \(\HiddenSpace{\sigma}\) with the same compact hidden body \(\HiddenBody{\sigma}\subset\HiddenSpace{\sigma}\), a diffeomorphism
    \[
    \ProtoProj{\sigma}\big|_{\ProtoZeroCore{\sigma}}
    :
    \ProtoZeroCore{\sigma}
    \xrightarrow{\cong}
    \GermZeroCore{\sigma},
    \]
    and a split diffeomorphism
    \[
    \ProtoSplit{\sigma}:
    \ProtoZeroCore{\sigma}\times\HiddenBody{\sigma}
    \xrightarrow{\cong}
    \ProtoShellSet{\sigma}
    \]
    obeying the projected split identity
    \[
    \GermSplit{\sigma}(z,\eta)
    =
    \ProtoProj{\sigma}\Bigl(
    \ProtoSplit{\sigma}\bigl(
    (\ProtoProj{\sigma}\big|_{\ProtoZeroCore{\sigma}})^{-1}(z),\eta
    \bigr)\Bigr)
    \]
    for every \(z\in\GermZeroCore{\sigma}\) and every \(\eta\in\HiddenBody{\sigma}\).
    \item A hidden charge
    \[
    \ProtoCharge{\sigma}:\HiddenSpace{\sigma}\to\EqnTarget{\sigma},
    \]
    a smooth observer-compatible exact proto-germ zero-response realization
    \[
    \ProtoEmbed{\sigma}:\ForSpace{\sigma}\hookrightarrow\ProtoZeroCore{\sigma},
    \]
    satisfying
    \[
    \GermCharge{\sigma}
    =
    \ProtoCharge{\sigma},
    \qquad
    \GermEmbed{\sigma}
    =
    \ProtoProj{\sigma}\circ\ProtoEmbed{\sigma},
    \]
    and a fixed-data solvability / coercive package on \(\ProtoZeroCore{\sigma}\), recorded through \(\ProtoKernelCh{\sigma}\) and \(\ProtoStateComp{\sigma}\), whose projection through
    \[
    \ProtoProj{\sigma}\big|_{\ProtoZeroCore{\sigma}}
    \]
    recovers the recorded germ kernel chart \(\GermKernelCh{\sigma}\) and orthogonal projector \(\GermStateComp{\sigma}\).
\end{enumerate}
The package is \emph{benchmark faithful} if it preserves the benchmark outputs \eqref{eq:fixedoutputs}, introduces no second operative metric, introduces no enlarged low-energy matter law, and is presented as a deeper substrate-side source for the current germ package rather than as a replacement benchmark theory.
\end{definition}

\begin{remark}
The label \emph{proto-germ} is used here only as a local marker for a deeper layer below the current germ frontier. It does not by itself assert a completed microscopic ontology or a final primitive law. Its role is only to name the next sharper transport package below the current germ dynamics.
\end{remark}

\section{Projected Exact-Shell Transport to the Germ Package}

\begin{proposition}[The germ shell is recovered by projected exact-shell transport]
\label{prop:shell}
Fix a benchmark-faithful proto-germ package. Then projected exact-shell transport along
\[
\ProtoProj{\sigma}\big|_{\ProtoShellSet{\sigma}}
\qquad\text{and}\qquad
\ProtoProj{\sigma}\big|_{\ProtoZeroCore{\sigma}}
\]
recovers the exact germ shell, the embedded germ zero core, and the germ split of Definition \ref{def:germ}.
\end{proposition}

\begin{proof}
Definition \ref{def:proto}(2) states that the restriction of \(\ProtoProj{\sigma}\) to \(\ProtoShellSet{\sigma}\) is a diffeomorphism onto \(\GermShellSet{\sigma}\). Definition \ref{def:proto}(3) states that the restriction of \(\ProtoProj{\sigma}\) to \(\ProtoZeroCore{\sigma}\) is a diffeomorphism onto \(\GermZeroCore{\sigma}\) and that the projected image of the proto-germ split is exactly the recorded germ split. These are precisely the claimed projected exact-shell transport statements.
\end{proof}

\begin{proposition}[Proto-germ dynamics canonically reproduce the current germ package]
\label{prop:germ}
Fix a benchmark-faithful proto-germ package. Then the projected exact-shell transport of Definition \ref{def:proto} reproduces exactly the current seed-generating substrate germ dynamics package of Definition \ref{def:germ}.
\end{proposition}

\begin{proof}
The germ body and the lower minimizer hierarchy are recovered by Definition \ref{def:proto}(1): the proto-germ body projects onto the recorded germ body, the proto-germ germ minimizer sections the projection on that body, and the projected proto-germ lower minimizers are exactly the recorded germ minimizers on the fixed seed, foundation, ground, origin, precursor, lift, and substrate fibers.

The germ restriction data are recovered by Definition \ref{def:proto}(2): the recorded germ section is the minimizing transport of the proto-germ section, the germ functional is therefore the projected restriction functional, and Proposition \ref{prop:shell} recovers the exact germ shell itself.

The remaining germ shell-internal data are recovered by Definition \ref{def:proto}(3)--(4): the embedded germ zero core is the projected proto-germ zero core, the germ split is the projected proto-germ split, the germ hidden charge is the proto-germ hidden charge, the exact germ zero-response realization is the projected proto-germ realization, and the fixed-data solvability / coercive package is the projection of the corresponding proto-germ package through the zero-core diffeomorphism. These are exactly the constituents retained by the current germ frontier.
\end{proof}

\begin{proposition}[The benchmark package remains fixed]
\label{prop:boundary}
For every benchmark-faithful proto-germ package, the operative metric remains exactly \(\CompletedMetric\), the ordinary matter route remains unchanged, no second operative metric is introduced, and the low-energy matter law is not enlarged.
\end{proposition}

\begin{proof}
This is part of the benchmark-faithfulness requirement in Definition \ref{def:proto}. Equation \eqref{eq:fixedoutputs} remains unchanged on the full primitive forcing shell, so the deeper proto-germ package does not alter the benchmark exact-closure package.
\end{proof}

\section{Main Theorem}

\begin{theorem}[Seed-generating substrate germ dynamics from germ-generating substrate proto-germ dynamics]
\label{thm:main}
Fix the primitive continuation data of Definition \ref{def:fixed}, the recorded lower-fiber ladder of Definition \ref{def:ladder}, the recorded downstream seed package of Definition \ref{def:seed}, and the recorded germ package of Definition \ref{def:germ}. Then, for each sector \(\sigma\in\{\Car,\Cur,\Hom\}\), every benchmark-faithful germ-generating substrate proto-germ dynamics package canonically produces the current germ package as projected exact-shell transport from the deeper proto-germ layer. More precisely:
\begin{enumerate}
    \item the current seed-generating substrate germ dynamics are reproduced unchanged, sector by sector, by projected exact-shell transport from the deeper proto-germ package;
    \item because the induced germ package is exactly the one already used by the current germ-frontier theorem, the current seed package, the minimal root exact-shell transport nucleus, and the previously recorded downstream root-side consequences remain available unchanged;
    \item the same benchmark one-metric package remains fixed: the operative metric is still exactly \(\CompletedMetric\), the ordinary matter route is unchanged, there is no second operative metric, and the low-energy matter law is not enlarged;
    \item consequently the remaining honest burden on the active line is deeper again: derive or justify the germ-generating substrate proto-germ dynamics themselves from still more primitive substrate structure, or else prove a genuinely smaller theorem-level collapse at that proto-germ level.
\end{enumerate}
\end{theorem}

\begin{proof}
Item (1) is Proposition \ref{prop:germ}. Item (2) is the routing consequence of item (1): the present note reproduces exactly the germ package that the current germ-frontier theorem already uses as its input, so that theorem and its downstream recovery of the seed package and the minimal root exact-shell transport nucleus apply unchanged. Item (3) is Proposition \ref{prop:boundary}. Item (4) is the project-level consequence of the previous items.
\end{proof}

\begin{corollary}[What remains open after this theorem]
\label{cor:next}
After Theorem \ref{thm:main}, the remaining honest burden is no longer to justify the seed-generating substrate germ dynamics themselves. Those dynamics are now the canonical projected exact-shell output of a sharper proto-germ class. What remains open is why the germ-generating substrate proto-germ dynamics should hold, or whether an even sharper theorem can remove that still deeper proto-germ-level freedom while preserving the same benchmark one-metric package and claim boundary.
\end{corollary}

\begin{proof}
Theorem \ref{thm:main} removes the current germ package as the unexplained stopping point on the active line. The next honest burden therefore lies strictly below it, unless a smaller theorem-level core inside the proto-germ package can first be isolated and proved sufficient.
\end{proof}

\section{Conclusion}

The positive result of this note is narrower than a completed first-principles derivation and stronger than another same-output germ restatement. The current germ-frontier theorem had already shown that the present seed package is the exact-shell transport image of a more primitive germ package. The present theorem then removes the next unexplained stopping point by showing that this germ package itself is the projected exact-shell transport image of a more primitive germ-generating substrate proto-germ dynamics.

The scope remains disciplined. The benchmark package is unchanged: the one operative metric is still exactly \(\CompletedMetric\), the ordinary matter route is unchanged, there is no second operative metric, and the low-energy matter law is not enlarged. Nothing here upgrades adoption to completed derivation or licenses stronger promotion language.

The open burden therefore moves one honest step further again. The next benchmark-facing manuscript should derive or justify the proto-germ dynamics from still more primitive substrate structure, unless the sharper honest gain available instead is a genuinely smaller theorem-level collapse inside the observer-relevant proto-germ package. Until then, the present result should be read for what it is: a deeper germ-scope derivation on the active line of \TheoryName{}, not a claim that the full first-principles program is finished.

\end{document}
